%% ============================================================
%%  Attention, Perception, & Psychophysics – Methods section
%%  Compiles out-of-the-box on Overleaf (no extra .cls needed).
%%  Uses APA 7th-edition author-year citations via natbib.
%%
%%  HOW TO USE IN OVERLEAF:
%%  1. Upload this file and refs.bib to your Overleaf project.
%%  2. Compile with pdfLaTeX.
%%  3. When ready to submit, paste the Method section into the
%%     official Springer Nature template (search "Springer Nature
%%     LaTeX Template" in the Overleaf gallery).
%% ============================================================

\documentclass[12pt, a4paper]{article}

\usepackage[margin=2.5cm]{geometry}
\usepackage{setspace}
\usepackage{times}          % Times New Roman (Springer preference)
\usepackage{microtype}
\usepackage{booktabs}
\usepackage{amsmath}
\usepackage[hidelinks]{hyperref}

%% APA author-year citations
\usepackage[
    natbib=true,
    style=apa,
    backend=biber
]{biblatex}
\addbibresource{refs.bib}

\doublespacing

%% ============================================================

\begin{document}

%% ---- TITLE PAGE (fill in before submission) ---------------
\title{Inhibition of return across six horizontal locations:\\
       A uniform spatial distribution}          % TODO: revise

\author{
    TODO First TODO Last\\
    \small TODO Department, TODO University
}
\date{}
\maketitle
\newpage

%% ============================================================
%%  ABSTRACT
%% ============================================================
\begin{abstract}
\noindent
TODO (150--250 words).
Inhibition of return (IOR)—the slowing of responses to targets at
previously cued locations—has been well established for single
peripheral locations.
Whether IOR is spatially uniform across multiple simultaneous
positions remains less clear.
The present study examined IOR across six equally spaced horizontal
locations (three per hemifield) using a detection task and a 300-ms
stimulus onset asynchrony (SOA).
[TODO: brief method summary, result summary, conclusion.]

\bigskip
\noindent\textit{Keywords:}
inhibition of return;
spatial attention;
detection task;
multiple locations;
exogenous cuing
\end{abstract}

\newpage

%% ============================================================
%%  METHOD
%% ============================================================

\section{Method}
\label{sec:method}

\subsection{Participants}

TODO participants (TODO female, TODO male;
$M_\mathrm{age} = \text{TODO}$~years, $SD = \text{TODO}$,
range TODO--TODO~years) were recruited TODO [e.g., from the departmental
participant pool / via Prolific].
All participants reported normal or corrected-to-normal vision and were
na\"{i}ve to the purpose of the study.
Informed consent was obtained prior to participation in accordance
with the Declaration of Helsinki, and the study was approved by the
TODO institutional ethics review board (approval no.\ TODO).

\subsection{Apparatus and Stimuli}

The experiment was programmed in JavaScript using jsPsych 7.3.4
\parencite{deleeuw2015jspsych} and presented in a web browser running
in fullscreen mode on participants' own computers.
Viewing distance was not physically constrained; participants were
instructed to maintain a comfortable, stable viewing position throughout
the session.

Stimuli were presented on a black background within a
$1000 \times 400$~pixel display area centred on the screen.
Six white-outlined placeholder boxes (each $80 \times 80$~px, 3-px
border) were arranged in a horizontal row.
Box centres were located at 100, 240, 380, 620, 760, and 900~px from
the left edge of the display area, labelled locations~1--6 from left
to right (see Figure~TODO).
The three leftward boxes (locations~1--3) and three rightward boxes
(locations~4--6) were separated by a 160-px edge-to-edge gap at the
vertical midline, which accommodated a white fixation cross
(\texttt{+}; 48-pt Arial) centred on the display.
Adjacent box centres within each hemifield were separated by 140~px.

The cue and target stimuli were identical white-filled squares
($74 \times 74$~px) rendered inside one of the six placeholder boxes,
filling the box interior while leaving the 3-px border visible.

\subsection{Design}

A $2 \times 6$ within-subjects design was employed with
\textit{cue validity} (valid, invalid) and \textit{location}
(1--6) as within-subjects factors.
On valid trials the target appeared at the cued location; on invalid
trials the target appeared at one of the five remaining locations,
selected at random.
A single cue-to-target stimulus onset asynchrony (SOA) of 300~ms was
used, consistent with the temporal window typically associated with
inhibition of return
\parencite[e.g.,][]{posner1984components, klein2000inhibition}.
The $2 \times 6 = 12$ unique conditions were each repeated 25 times,
yielding 300 experimental trials per participant.
An additional 60 catch trials (20\,\% of the experimental trial count)
were intermixed uniformly throughout the session to discourage
anticipatory responding (see Procedure section).

\subsection{Procedure}
\label{sec:procedure}

After providing informed consent and entering fullscreen mode,
participants read on-screen instructions describing the task.
They then completed 12 practice trials (data not analysed) to
familiarise themselves with the trial sequence.
The 360 main trials (300 experimental, 60 catch) were presented in a
single continuous block in randomised order.
Self-paced rest breaks were offered automatically every 48~trials.

Each trial proceeded as follows (see Figure~TODO):
\begin{enumerate}
    \item \textit{Fixation} (750~ms): the fixation cross and six empty
          placeholder boxes were displayed.
    \item \textit{Cue} (50~ms): a white square appeared inside one of
          the six boxes.
    \item \textit{Cue--target interval} (250~ms): the display returned
          to the fixation cross and empty boxes.
          Together with the cue duration, this produced a cue-to-target
          SOA of 300~ms.
    \item \textit{Target/catch period}: on experimental trials, a white
          square appeared inside the target box and remained until the
          participant responded (no deadline imposed).
          On catch trials, the fixation display continued for 1000~ms
          with no target appearing.
    \item \textit{Inter-trial interval} (500~ms): blank black screen.
\end{enumerate}

Participants were instructed to keep their gaze on the fixation cross
at all times, to ignore the cue, and to press the spacebar as quickly
as possible upon detecting the target.
They were explicitly informed that on a minority of trials no target
would appear, and that they should withhold their response on those
trials.

\subsection{Data Analysis}

Trials on which the reaction time (RT) was less than 100~ms were
classified as anticipatory and excluded from all RT analyses.
Catch-trial false alarms (spacebar presses during the 1000-ms
no-target window) were tallied per participant and are reported
descriptively as a measure of response caution.

Mean correct RTs were submitted to a $2 \times 6$
repeated-measures analysis of variance (ANOVA) with
\textit{validity} (valid, invalid) and \textit{location} (1--6)
as within-subjects factors.
Inhibition of return was operationalised as reliably slower RTs on
valid (cued) relative to invalid (uncued) trials at the 300-ms SOA
\parencite{posner1984components}.
The \textit{Validity} $\times$ \textit{Location} interaction was used
to assess whether IOR magnitude varied systematically across spatial
positions; a null interaction was taken as evidence of spatially
uniform IOR.
Effect sizes are reported as partial eta-squared
($\eta^{2}_{p}$).
All analyses were conducted in TODO [e.g.,
\textit{R}~4.x \parencite{rcore2024} /
SPSS~29 / Python~3.x with \texttt{pingouin}].

%% ============================================================
%%  REFERENCES
%% ============================================================

\printbibliography

\end{document}
